% ID: SEL-D0064-OA-E02
% ESTADO: cotejada. Segunda pasada visual y borrador compilado revisados.
% Fuente: Selectividad/sel_2010_matematicas_aplicadas/mat_apli/mat_apli_sept/mat_apli_sept/EXAMEN-ANDALUCÍA-3.pdf
% SHA-256: 9c3d2d628115a2f20f17b47171106ddd90f214b5f6804bdf9a7dcb7037c75efd
% Página: 1; opción A; ejercicio original 2.

Un consultorio médico abre a las 5 de la tarde y cierra cuando no hay pacientes.

La expresión que representa el número medio de pacientes en función del tiempo en horas, $t$, que lleva abierto el consultorio es $N(t)=4t-t^2$
\begin{enumerate}[label=\alph*)]
\item \textbf{(1 punto)} \textquestiondown{}A qué hora el número medio de pacientes es máximo? \textquestiondown{}Cuál es ese máximo?
\item \textbf{(1 punto)} Sabiendo que el consultorio cierra cuando no hay pacientes, \textquestiondown{}a qué hora cerrará?
\item \textbf{(0.5 puntos)} Represente gráficamente $N(t)=4t-t^2$, con $N(t)\geq0$.
\end{enumerate}
